\section{Подготовка к работе}

Перед тем как приступить к написанию кода и практическим полетам, необходимо тщательно подготовить рабочее место и изучить особенности эксплуатации оборудования. В этом разделе мы подробно разберем правила техники безопасности, установку необходимого программного обеспечения и процесс загрузки программ на борт квадрокоптера.

\subsection{Техника безопасности}

\begin{warnbox}[Критически важно!]
    Соблюдение правил безопасности — это фундаментальный навык любого инженера и пилота БПЛА. Пренебрежение этими правилами может привести к травмам или поломке дорогостоящего оборудования.
\end{warnbox}

Квадрокоптер «Геоскан Пионер» является сложным электромеханическим устройством. При работе с ним необходимо неукоснительно соблюдать следующие требования:

\begin{itemize}
    \item \textbf{Работа с механикой только без питания:} Любые монтажные работы, замена деталей или подключение модулей расширения должны производиться на полностью обесточенном аппарате (аккумулятор отключен).
    
    \item \textbf{Снятие пропеллеров при отладке:} Это «золотое правило» программиста дронов.
    \begin{itemize}
        \item При отладке скриптов, тестировании моторов или проверке реакции на команды пульта \textbf{пропеллеры должны быть сняты}.
        \item Ремонт, сборка и разборка квадрокоптера всегда осуществляются со снятыми пропеллерами.
        \item Случайный запуск двигателей (например, из-за ошибки в коде или случайного касания стика газа) при установленных пропеллерах неминуемо приведет к неконтролируемому взлету, травмам пальцев или повреждению имущества.
    \end{itemize}

    \item \textbf{Обращение с аккумуляторами:} Используйте только исправные литий-полимерные (LiPo) батареи. Не оставляйте их на зарядке без присмотра. Если аккумулятор вздулся или имеет механические повреждения оболочки — его эксплуатация строго запрещена.

    \item \textbf{Порядок включения и выключения:}
    \begin{itemize}
        \item \textbf{Включение:} Сначала всегда включается пульт управления (передатчик), и только затем подается питание на борт дрона. Это гарантирует, что дрон сразу увидит управляющий сигнал.
        \item \textbf{Выключение:} В обратном порядке — сначала отключаем питание дрона, затем выключаем пульт.
    \end{itemize}

    \item \textbf{Организация рабочего пространства:} Уберите со стола лишние предметы, провода и инструменты. Дрон, даже без пропеллеров, может создавать вибрацию, а с пропеллерами — сильный воздушный поток, который может сдуть мелкие детали или затянуть легкие предметы в винты.
\end{itemize}

\subsection{Инструментарий разработчика}

Для успешного освоения курса вам потребуется специализированное программное обеспечение. Мы разделяем инструменты на те, что необходимы для работы непосредственно с «Пионером», и те, что удобны для изучения чистого языка Lua.

\subsubsection{Работа с квадрокоптером: Pioneer Station}

Основной средой разработки для квадрокоптера «Геоскан Пионер» является \textbf{Pioneer Station}. Это интегрированная среда (IDE), предоставляющая полный цикл инструментов для взаимодействия с автопилотом:

\begin{itemize}
    \item \textbf{Редактор кода:} Позволяет писать скрипты на языке Lua с подсветкой синтаксиса.
    \item \textbf{Загрузчик:} Обеспечивает компиляцию и передачу программ в память дрона.
    \item \textbf{Терминал и отладка:} Выводит сообщения из функции \texttt{print()} и позволяет следить за состоянием датчиков (гироскопа, акселерометра, дальномера) в реальном времени.
    \item \textbf{Менеджер прошивок:} Позволяет обновлять встроенное ПО автопилота до актуальной версии.
\end{itemize}

\begin{infobox}[Установка]
    Убедитесь, что на вашем компьютере установлена последняя версия Pioneer Station. Дистрибутив можно бесплатно скачать с официального сайта поддержки Геоскан.
\end{infobox}

\subsubsection{Изучение синтаксиса Lua}

Язык Lua является стандартным, и его синтаксис (циклы, условия, функции) можно изучать отдельно от квадрокоптера. Это часто бывает удобнее, так как не требует подключения аккумулятора и постоянной перезагрузки дрона.

Для тренировки в написании чистого кода на Lua (без специфичных функций дрона, таких как \texttt{ap.goToLocalPoint}) мы рекомендуем использовать:

\begin{itemize}
    \item \textbf{ZeroBrane Studio:} Легковесная и быстрая IDE, созданная специально для Lua. Отлично подходит для новичков.
    \item \textbf{Visual Studio Code:} Популярный редактор кода. Для работы с Lua потребуется установить расширение (например, \textit{Lua by sumneko}).
    \item \textbf{Онлайн-компиляторы:} Для быстрой проверки коротких фрагментов кода можно использовать веб-сервисы, такие как \textit{Lua Demo} или \textit{Replit}.
\end{itemize}

\textit{Примечание: Код, написанный в сторонних редакторах, для запуска на дроне все равно необходимо будет перенести в Pioneer Station.}

\subsection{Подключение квадрокоптера}

Процесс соединения дрона с компьютером состоит из нескольких простых шагов:

\begin{enumerate}
    \item Убедитесь, что пропеллеры сняты (если вы не планируете немедленный вылет на полигоне).
    \item Подключите аккумулятор к разъему питания квадрокоптера. Вы услышите характерную мелодию инициализации регуляторов оборотов, и светодиодная индикация начнет мигать.
    \item Подключите квадрокоптер к USB-порту компьютера с помощью кабеля Micro-USB.
    \item Запустите приложение Pioneer Station.
    \item В интерфейсе программы выберите соответствующий COM-порт (обычно он определяется автоматически при подключении кабеля).
\end{enumerate}

\subsection{Ваша первая программа}

Чтобы убедиться, что все системы работают исправно, напишем и запустим простейшую программу. Традиционно в программировании первой программой является вывод приветствия.

Алгоритм запуска скрипта:
\begin{enumerate}
    \item \textbf{Написание:} Откройте редактор в Pioneer Station и введите код программы.
    \item \textbf{Загрузка:} Нажмите кнопку «Загрузить» (иконка Play или стрелка). Скрипт будет передан в память автопилота.
    \item \textbf{Исполнение:} После загрузки дрон автоматически перезагрузит виртуальную машину Lua и начнет выполнение скрипта. Результат работы вы увидите в окне «Консоль» (Log).
\end{enumerate}

\begin{codebox}[Пример: Hello World]
\begin{codefile}{lua/examples/chapter_02/getting_started_hello.lua}
\end{codefile}
\end{codebox}

Если в логе появилось ваше сообщение — поздравляем! Вы успешно настроили рабочее место и готовы погрузиться в мир программирования полетных контроллеров.
